\documentclass[9pt,aspectratio=169,xcolor=table]{beamer}
\usepackage[compact]{styles/stmcommon}

\title{Designing, Building and Verifying an Embedded System}
\subtitle{\emph{Introduction to ARM Cortex-M Microprocessors} --- Chapter 14}

\begin{document}

\titlepage

%===============================================================================
\section{Why This Matters}
%===============================================================================

\begin{frame}{Why This Chapter Matters}
    \begin{columns}[T]
    \begin{column}{.47\textwidth}
        {\normalsize Earlier chapters used one peripheral at a time. A real system samples on a schedule, answers buttons, drives indicators and reports to a host, all at once, on one processor, with no operating system. The hard part is deciding \emph{before} coding which peripherals the job needs and whether the timing can work. Then you must \emph{prove} the result meets the requirement.}
    \end{column}
    \begin{column}{.51\textwidth}
        \begin{block}{After this chapter, you can:}
        {\small
        \begin{enumerate}\itemsep2pt
            \item Turn a requirement into a decomposition and a peripheral choice.
            \item Choose polling, timed polling or interrupts per activity.
            \item Structure a cooperative, non-blocking main loop.
            \item Protect data shared between ISR and main.
            \item Write a timing budget, then verify it by measurement.
            \item Diagnose a multi-peripheral system bottom-up.
        \end{enumerate}}
        \end{block}
    \end{column}
    \end{columns}
\end{frame}


%===============================================================================
\section{From Requirement to Architecture}
%===============================================================================

\begin{frame}{The Step Before Any Code}
    \leavevmode
    A requirement is prose. An architecture is a set of decisions. Getting from
    one to the other is \alert{the engineering}, and it happens on paper.

    \vspace{3mm}
    \begin{center}
    \begin{tikzpicture}[x=10mm, y=9mm]
        \node[blk,    minimum width=28mm] (r) at (0,1.5)  {Written requirement};
        \node[blk,    minimum width=28mm] (f) at (0,0.5)  {Functional decomposition};
        \node[blk,    minimum width=28mm] (p) at (0,-0.5) {Peripheral selection};
        \node[cpublk, minimum width=28mm] (s) at (0,-1.5) {Program structure};
        \foreach \x/\y in {r/f, f/p, p/s}{\draw[flow] (\x) -- (\y);}
        \node[lbl, anchor=west] at (1.65,1.0)  {what activities exist?};
        \node[lbl, anchor=west] at (1.65,0.0)  {what hardware does each need?};
        \node[lbl, anchor=west] at (1.65,-1.0) {how are they scheduled?};
    \end{tikzpicture}
    \end{center}

    \vspace{2mm}
    \begin{block}{Why this order matters}
        Choosing peripherals before decomposing the activities means you discover
        the pin conflict \alert{after} the board is made.
    \end{block}
\end{frame}

\begin{frame}{Four Ways to Structure an Activity}
    \begin{center}
    {\small
    \striped
    \begin{tabular}{@{}l R{.34\textwidth} R{.32\textwidth}@{}}
        \toprule
        \textbf{Structure} & \textbf{Use when} & \textbf{Cost} \\
        \midrule
        \textbf{Blocking poll} & Never, in a real system & Burns the CPU; blocks everything else \\
        \textbf{Timed poll} & The activity has a known period and loose jitter & Simple; jitter is the loop time \\
        \textbf{Interrupt} & The event is external and the deadline is tight & Needs care with shared data \\
        \textbf{DMA} & Bulk data, no per-byte decision & Setup cost; fewest CPU cycles \\
        \bottomrule
    \end{tabular}}
    \end{center}

    \vspace{3mm}
    \begin{block}{The decision is per activity, not per program}
        One system routinely uses all four: DMA for an ADC stream, an interrupt
        for a button, a timed poll for a display refresh, and the main loop for
        everything with no deadline at all.
    \end{block}
\end{frame}


%===============================================================================
\section{The Project: A Sensor Monitor}
%===============================================================================

\begin{frame}{The Project}
    \begin{center}
    \begin{tikzpicture}[x=10mm, y=10mm]
        \node[cpublk, minimum width=36mm] (mcu) at (0,2) {STM32 Cortex-M4};
        \node[blk, minimum width=18mm] (adc)   at (-5.2,0) {ADC};
        \node[blk, minimum width=18mm] (tim)   at (-2.6,0) {TIM2};
        \node[blk, minimum width=18mm] (gpio)  at ( 0.0,0) {GPIO};
        \node[blk, minimum width=18mm] (usart) at ( 2.6,0) {USART1};
        \node[blk, minimum width=18mm] (exti)  at ( 5.2,0) {EXTI};
        \coordinate (bus) at ($(mcu.south)+(0,-0.6)$);
        \draw[busline] (mcu.south) -- (bus);
        \draw[busline] (adc.north |- bus) -- (exti.north |- bus);
        \foreach \p in {adc,tim,gpio,usart,exti}{\draw[flow] (\p.north |- bus) -- (\p.north);}
        \node[lbl, below=4mm of adc,   align=center] {potentiometer\\or sensor};
        \node[lbl, below=4mm of tim,   align=center, text=accentline] {1~kHz tick\\{\scriptsize (internal)}};
        \node[lbl, below=4mm of gpio,  align=center] {4 LED\\bar};
        \node[lbl, below=4mm of usart, align=center] {PC terminal\\9600 8N1};
        \node[lbl, below=4mm of exti,  align=center] {button\\PC13};
        \foreach \p in {adc,gpio,usart,exti}{\draw[flow] (\p.south) -- ++(0,-0.5);}
        \draw[flow, draw=accentline, dashed] (tim.south) -- ++(0,-0.5);
    \end{tikzpicture}
    \end{center}

    \vspace{2mm}
    {\small \textbf{Five peripherals, one processor, no OS.}}

    \vspace{2mm}
    {\footnotesize
    \begin{enumerate}\itemsep1pt
        \item Sample the ADC at exactly \alert{10~Hz}.
        \item Display the level as a 4-LED bar.
        \item Transmit each reading over UART as text.
        \item Button press toggles report between \alert{raw counts} and \alert{millivolts}.
        \item Never miss a sample. Never block.
    \end{enumerate}}
\end{frame}


%===============================================================================
\section{Architecture: The Non-Blocking Main Loop}
%===============================================================================

\begin{frame}[fragile]{The Pattern}
\begin{lstlisting}[language=C, basicstyle=\ttfamily\scriptsize]
while (1) {
    uint32_t now = millis();

    if (now - last_sensor >= 2000) {    // every 2 s
        last_sensor = now;
        read_sensor();
    }
    if (now - last_display >= 100) {    // every 100 ms
        last_display = now;
        update_display();
    }
    if (button_flag) {                  // set by an ISR
        button_flag = 0;
        handle_button();
    }
}
\end{lstlisting}

    \vspace{2mm}
    \begin{block}{Nothing waits}
        Each activity asks ``is it my turn yet?'' and returns immediately if not.
        The loop runs thousands of times per second, and every activity gets
        serviced \alert{close enough} to its deadline.
    \end{block}
\end{frame}

\begin{frame}{Why \texttt{now - last >= period} and Not \texttt{now >= next}}
    \begin{columns}[T]
    \begin{column}{.48\textwidth}
        \textbf{The subtraction form}

        \vspace{1mm}
        {\ttfamily\scriptsize if (now - last >= 2000)}

        \vspace{2mm}
        {\small Survives the 32-bit counter \alert{wrapping} after 49 days,
        because unsigned subtraction wraps consistently.}
    \end{column}
    \begin{column}{.48\textwidth}
        \textbf{The comparison form}

        \vspace{1mm}
        {\ttfamily\scriptsize if (now >= next)}

        \vspace{2mm}
        {\small \alert{Breaks} at the wrap: \texttt{next} is a huge number,
        \texttt{now} has reset to a small one, and the activity stops for
        49 days.}
    \end{column}
    \end{columns}

    \vspace{3mm}
    \begin{block}{A bug you cannot find by testing}
        It appears once, after seven weeks of continuous running. This is exactly
        the class of fault that \alert{reasoning} catches and \alert{testing} does
        not --- which is why the argument matters more than the code.
    \end{block}
\end{frame}

\begin{frame}{Sharing Data With an ISR}
    \begin{columns}[T]
    \begin{column}{.48\textwidth}
        \textbf{Two hazards}
        {\small
        \begin{itemize}\itemsep1pt
            \item \alert{Missing \texttt{volatile}} --- the optimiser caches a
                  value the ISR changes
            \item \alert{Torn access} --- \texttt{main()} reads a multi-word
                  variable while the ISR is halfway through writing it
        \end{itemize}}
    \end{column}
    \begin{column}{.48\textwidth}
        \textbf{Three defences}
        {\small
        \begin{itemize}\itemsep1pt
            \item Mark it \texttt{volatile}
            \item Keep it to a single \texttt{uint32\_t} --- read atomically on a
                  32-bit machine
            \item Or disable interrupts \alert{briefly} around the access
        \end{itemize}}
    \end{column}
    \end{columns}

    \vspace{3mm}
    \begin{block}{The flag pattern}
        The ISR does the minimum --- clear the hardware flag, set a software one,
        return. \texttt{main()} does the real work. That keeps the handler short
        \alert{and} keeps the shared state to one byte.
    \end{block}
\end{frame}


%===============================================================================
\section{The Complete Program}
%===============================================================================

\begin{frame}[fragile]{The Complete Program (1)}{Shared state and interrupt handlers}
    \begin{columns}[T]
    \begin{column}{.5\textwidth}
\begin{cc}[basicstyle=\ttfamily\scriptsize]
#include "stm32f4xx.h"

/* Shared state: written by ISRs, read by main */
volatile uint32_t msTicks     = 0;
volatile uint16_t adc_result  = 0;
volatile uint8_t  adc_ready   = 0;
volatile uint8_t  button_flag = 0;

/* Application state (main context only) */
static uint8_t show_millivolts = 0;

void SysTick_Handler(void) {
    msTicks++;      // its only job: advance time
}

void ADC_IRQHandler(void) {
    if (ADC1->SR & (1U << 1)) {   // EOC
        adc_result = (uint16_t)ADC1->DR;
        adc_ready  = 1;
    }
}
\end{cc}
    \end{column}
    \begin{column}{.46\textwidth}
\begin{cc}[basicstyle=\ttfamily\scriptsize]
void EXTI15_10_IRQHandler(void) {
    static uint32_t last_press = 0;

    if (EXTI->PR & (1U << 13)) {
        EXTI->PR = (1U << 13);  // clear FIRST

        // Debounce by timestamp,
        // no delay loop in an ISR.
        if ((msTicks - last_press) > 200) {
            last_press  = msTicks;
            button_flag = 1;
        }
    }
}
\end{cc}
        \vspace{1mm}
        {\small Each handler does the \alert{minimum}: capture the event, set a \texttt{volatile} flag, return.}
    \end{column}
    \end{columns}
\end{frame}

\begin{frame}[fragile]{The Complete Program (2)}{Setup: LEDs and button}
    \begin{columns}[T]
    \begin{column}{.46\textwidth}
\begin{cc}[basicstyle=\ttfamily\scriptsize]
static void leds_init(void) {
    RCC->AHB1ENR |= (1U << 1);   // GPIOB
    for (uint32_t pin = 0; pin < 4; pin++) {
        GPIOB->MODER &= ~(3U << (pin * 2));
        GPIOB->MODER |=  (1U << (pin * 2));
    }
}
\end{cc}
        \vspace{1mm}
        {\small Four outputs, \texttt{PB0}--\texttt{PB3}: the LED bar.}
    \end{column}
    \begin{column}{.5\textwidth}
\begin{cc}[basicstyle=\ttfamily\scriptsize]
static void button_init(void) {
    RCC->AHB1ENR |= (1U << 2);   // GPIOC
    GPIOC->MODER &= ~(3U << (13 * 2));
    GPIOC->PUPDR &= ~(3U << (13 * 2));
    GPIOC->PUPDR |=  (1U << (13 * 2));
    RCC->APB2ENR |= (1U << 14);  // SYSCFG
    SYSCFG->EXTICR[3] &= ~(0xFU << 4);
    SYSCFG->EXTICR[3] |=  (2U << 4);
    EXTI->FTSR |= (1U << 13);    // falling
    EXTI->IMR  |= (1U << 13);    // unmask
    NVIC_SetPriority(EXTI15_10_IRQn, 3);
    NVIC_EnableIRQ(EXTI15_10_IRQn);
}
\end{cc}
    \end{column}
    \end{columns}
    \vspace{1mm}
    {\small Every peripheral follows the same shape: \alert{clock, pins, parameters, then use}.}
\end{frame}

\begin{frame}[fragile]{The Complete Program (3)}{Setup: ADC and UART}
    \begin{columns}[T]
    \begin{column}{.46\textwidth}
\begin{cc}[basicstyle=\ttfamily\scriptsize]
static void adc_init(void) {
    RCC->AHB1ENR |= (1U << 0);   // GPIOA
    RCC->APB2ENR |= (1U << 8);   // ADC1
    GPIOA->MODER |= (3U << (0 * 2));
    ADC1->SQR1  &= ~(0xFU << 20);
    ADC1->SQR3  &= ~(0x1FU << 0);
    ADC1->SMPR2 &= ~(7U << 0);
    ADC1->SMPR2 |=  (7U << 0);   // 480 cycles
    ADC1->CR1   |=  (1U << 5);   // EOCIE
    ADC1->CR2   |=  (1U << 0);   // ADON
    NVIC_SetPriority(ADC_IRQn, 1);
    NVIC_EnableIRQ(ADC_IRQn);
}
\end{cc}
        \vspace{1mm}
        {\small ADC priority 1, button priority 3: the ADC is more urgent.}
    \end{column}
    \begin{column}{.5\textwidth}
\begin{cc}[basicstyle=\ttfamily\scriptsize]
static void uart_init(void) {
    RCC->AHB1ENR |= (1U << 1);   // GPIOB
    RCC->APB2ENR |= (1U << 4);   // USART1
    GPIOB->MODER &= ~((3U<<(6*2)) | (3U<<(7*2)));
    GPIOB->MODER |=  ((2U<<(6*2)) | (2U<<(7*2)));
    GPIOB->AFR[0] &= ~((0xFU<<(6*4)) | (0xFU<<(7*4)));
    GPIOB->AFR[0] |=  ((7U<<(6*4)) | (7U<<(7*4)));
    USART1->BRR  = (104U << 4) | 3U; // 9600 @ 16 MHz
    USART1->CR1 |= (1U << 3) | (1U << 2); // TE, RE
    USART1->CR1 |= (1U << 13);   // UE
}
\end{cc}
        \vspace{1mm}
        {\small Pin map: \texttt{PA0} analog, \texttt{PB6}/\texttt{PB7} USART1, \texttt{PB0}--\texttt{PB3} LED bar, \texttt{PC13} button.}
    \end{column}
    \end{columns}
\end{frame}

\begin{frame}[fragile]{The Complete Program (4)}{Helpers}
    \begin{columns}[T]
    \begin{column}{.5\textwidth}
\begin{cc}[basicstyle=\ttfamily\scriptsize]
static void uart_write(uint8_t c) {
    while (!(USART1->SR & (1U << 7))) { }  // TXE
    USART1->DR = c;
}

static void uart_print(const char *s) {
    while (*s) uart_write((uint8_t)*s++);
}

static void uart_print_uint(uint32_t v) {
    char buf[11];
    int i = 0;
    if (v == 0) { uart_write('0'); return; }
    while (v) { buf[i++] = '0' + (v % 10); v /= 10; }
    while (i)  { uart_write((uint8_t)buf[--i]); }
}
\end{cc}
    \end{column}
    \begin{column}{.46\textwidth}
\begin{cc}[basicstyle=\ttfamily\scriptsize]
// 12-bit reading -> 4-LED bar,
// written atomically via BSRR
static void leds_show_bar(uint16_t raw) {
    uint32_t level = (raw * 5U) / 4096U;  // 0..4
    uint32_t pattern = (level >= 4) ? 0x0F
                       : ((1U << level) - 1U);
    GPIOB->BSRR = (pattern & 0x0F)
                | ((~pattern & 0x0F) << 16);
}
\end{cc}
        \vspace{1mm}
        {\small One \texttt{BSRR} store sets some LEDs and clears the others in the same cycle.}
    \end{column}
    \end{columns}
\end{frame}

\begin{frame}[fragile]{The Complete Program (5)}{\texttt{main()}: setup and two of the three tasks}
    \begin{columns}[T]
    \begin{column}{.6\textwidth}
\begin{cc}[basicstyle=\ttfamily\scriptsize]
int main(void) {
    __disable_irq();
    leds_init(); button_init();
    adc_init();  uart_init();
    SysTick_Config(16000000 / 1000);  // 1 ms tick
    __enable_irq();
    uart_print("\r\nSensor Monitor ready.\r\n");

    uint32_t next_sample = 0;
    while (1) {
        /* Task 1: start a conversion every 100 ms */
        if ((msTicks - next_sample) >= 100) {
            next_sample += 100;       // no drift
            ADC1->CR2 |= (1U << 30);  // SWSTART
        }

        /* Task 2: a result arrived (next slide) */

        /* Task 3: button pressed */
        if (button_flag) {
            button_flag = 0;
            show_millivolts = !show_millivolts;
            uart_print(show_millivolts
                       ? "-> millivolts\r\n" : "-> raw\r\n");
        }
    }
}
\end{cc}
    \end{column}
    \begin{column}{.36\textwidth}
        \textbf{Read it as three \texttt{if}s}
        {\small
        \begin{itemize}\itemsep2pt
            \item \texttt{(msTicks - next\_sample) >= 100}: a sample is \alert{due}.
            \item \texttt{adc\_ready}: a result \alert{arrived}.
            \item \texttt{button\_flag}: the user \alert{pressed}.
        \end{itemize}}
        \vspace{1mm}
        {\small Each branch finishes quickly and returns to the top. \alert{No branch waits for another.}}
    \end{column}
    \end{columns}
\end{frame}

\begin{frame}[fragile]{The Complete Program (6)}{Task 2: a result arrived}
    \begin{columns}[T]
    \begin{column}{.58\textwidth}
\begin{cc}[basicstyle=\ttfamily\scriptsize]
if (adc_ready) {
    uint16_t raw;

    __disable_irq();            // critical section:
    raw       = adc_result;     // copy the shared value
    adc_ready = 0;
    __enable_irq();

    leds_show_bar(raw);

    if (show_millivolts) {
        // Integer maths only: no FPU needed.
        uint32_t mv = ((uint32_t)raw * 3300U) / 4095U;
        uart_print("Level: ");
        uart_print_uint(mv);
        uart_print(" mV\r\n");
    } else {
        uart_print("Raw:   ");
        uart_print_uint(raw);
        uart_print("\r\n");
    }
}
\end{cc}
    \end{column}
    \begin{column}{.38\textwidth}
        \begin{block}{Why a critical section?}
            {\small The ISR may fire between reading the value and clearing the flag. Interrupts are off for only \alert{three statements}. The slow UART printing happens outside.}
        \end{block}
    \end{column}
    \end{columns}
\end{frame}


%===============================================================================
\section{Design Decisions Worth Noting}
%===============================================================================

\begin{frame}{Design Decisions Worth Noting}
    \begin{center}
    {\footnotesize
    \striped
    \begin{tabular}{@{}l R{.7\textwidth}@{}}
        \toprule
        \textbf{Decision} & \textbf{Why} \\
        \midrule
        \textbf{The schedule does not drift} & \texttt{next\_sample += 100} steps by a fixed amount. \texttt{msTicks + 100} would add any delay from the rest of the loop, and the rate would slowly fall below 10\,Hz. That breaks the uniform interval Nyquist relies on. \\
        \textbf{A minimal critical section} & \texttt{adc\_result} and \texttt{adc\_ready} are read and cleared with interrupts off, so the ISR cannot overwrite the value between the two accesses. Only three statements are protected. UART printing happens outside. \\
        \textbf{Priorities follow urgency} & ADC priority 1, button priority 3. A conversion result is lost if not read before the next one. A button press can wait. \\
        \textbf{Integer maths, no FPU} & \texttt{(raw * 3300) / 4095}: multiply first to keep precision. $4095\times3300$ fits in 32 bits. Dividing first would truncate to zero. \\
        \textbf{The LED bar is atomic} & One \texttt{BSRR} store updates all four LEDs, so no half-finished pattern is ever visible. \\
        \bottomrule
    \end{tabular}}
    \end{center}
\end{frame}

\begin{frame}{You Already Know This: A State Machine}
    \leavevmode
    This architecture is a \alert{state machine with a shared time base}. You have designed exactly this, on paper, as a synchronous sequential circuit.

    \vspace{3mm}
    \begin{center}
    {\small
    \striped
    \begin{tabular}{@{}l l@{}}
        \toprule
        \textbf{In digital electronics} & \textbf{In this program} \\
        \midrule
        Edge-triggered inputs & Interrupts that set flags \\
        Combinational next-state logic & The \texttt{if} tests in the main loop \\
        The clock & SysTick \\
        State in flip-flops & State in \texttt{volatile} variables \\
        Change behaviour: rewire & Change behaviour: \alert{edit the code} \\
        \bottomrule
    \end{tabular}}
    \end{center}
\end{frame}


%===============================================================================
\section{The Timing Budget}
%===============================================================================

\begin{frame}{The Timing Budget}
    \leavevmode
    Write it \alert{before} implementation, and it tells you whether the design is
    possible at all.

    \vspace{3mm}
    \begin{center}
    {\small
    \striped
    \begin{tabular}{@{}l l l l@{}}
        \toprule
        \textbf{Activity} & \textbf{Period} & \textbf{Duration} & \textbf{Load} \\
        \midrule
        SysTick ISR      & 1\,ms   & 2\,\textmu s   & 0.2\,\% \\
        Sensor read      & 2\,s    & 400\,\textmu s & 0.02\,\% \\
        Display update   & 100\,ms & 3\,ms          & 3\,\% \\
        Button ISR       & rare    & 5\,\textmu s   & $\approx 0$ \\
        \midrule
        \textbf{Total}   &         &                & \alert{$\approx 3.3\,\%$} \\
        \bottomrule
    \end{tabular}}
    \end{center}

    \vspace{3mm}
    \begin{block}{What the number tells you}
        3\,\% means comfortable headroom. \alert{80\,\% means the design is
        fragile} --- one slow path and a deadline is missed. Over 100\,\% means it
        cannot work, and you know that before writing a line.
    \end{block}
\end{frame}


%===============================================================================
\section{Predict, Then Measure}
%===============================================================================

\begin{frame}{Predict, Then Measure}
    \begin{columns}[T]
    \begin{column}{.48\textwidth}
        \textbf{The discipline}
        \begin{enumerate}
            \item \alert{Predict} the number --- period, duty cycle, execution
                  time
            \item \alert{Measure} it
            \item When they disagree, \alert{find out why}
        \end{enumerate}
    \end{column}
    \begin{column}{.48\textwidth}
        \textbf{How to measure}
        {\small
        \begin{itemize}\itemsep1pt
            \item Set a GPIO pin high on entry, low on exit --- watch it on a
                  scope
            \item Read \texttt{SysTick->VAL} before and after
            \item Toggle a pin per loop iteration and measure the frequency
        \end{itemize}}
    \end{column}
    \end{columns}

    \vspace{3mm}
    \begin{block}{The disagreement is the finding}
        If you predicted 400\,\textmu s and measure 3\,ms, you have learned
        something --- about the compiler, the clock tree, or your own assumptions.

        \vspace{1mm}
        \alert{Measuring without predicting teaches you nothing}, because any
        number looks plausible.
    \end{block}
\end{frame}


%===============================================================================
\section{What a Working Demonstration Does Not Prove}
%===============================================================================

\begin{frame}{What a Working Demonstration Does Not Prove}
    \begin{columns}[T]
    \begin{column}{.5\textwidth}
        \textbf{It runs. What is still unknown?}
        {\small
        \begin{itemize}\itemsep1pt
            \item Whether it survives the \alert{counter wrap} at 49 days
            \item Whether the deadline is met in the \alert{worst case}, not the
                  typical one
            \item Whether a shared variable is ever \alert{torn}
            \item What happens when a sensor \alert{fails} rather than reads wrong
            \item Whether it works at the \alert{temperature extremes}
        \end{itemize}}
    \end{column}
    \begin{column}{.46\textwidth}
        \begin{block}{Demonstration is evidence, not proof}
            A demo shows the system works \alert{once, here, now}, under the
            conditions you happened to create.

            \vspace{2mm}
            Everything else you believe about it comes from \alert{argument}: the
            timing budget, the wrap analysis, the reasoning about shared state.
        \end{block}
    \end{column}
    \end{columns}
\end{frame}


%===============================================================================
\section{A Debugging Methodology}
%===============================================================================

\begin{frame}{A Debugging Methodology}
    \leavevmode
    When a multi-peripheral system misbehaves, work \alert{bottom-up}. Each step
    either finds the fault or eliminates a layer.

    \vspace{1mm}
    \begin{center}
    {\small
    \striped
    \begin{tabular}{@{}l R{.60\textwidth}@{}}
        \toprule
        \textbf{Step} & \textbf{Question} \\
        \midrule
        1 & Is the \alert{clock} enabled for this peripheral? \\
        2 & Is the \alert{pin} in the right mode, with the right alternate function? \\
        3 & Does the \alert{register} read back what you wrote? \\
        4 & Does the peripheral's \alert{status flag} ever set? \\
        5 & Does the \alert{ISR} run --- toggle a pin inside it and look \\
        6 & Is the \alert{shared variable} \texttt{volatile}, and read atomically? \\
        7 & Does the \alert{timing} match the budget you predicted? \\
        \bottomrule
    \end{tabular}}
    \end{center}

    \vspace{1mm}
    {\small Steps 1 and 2 account for most first-time faults --- and
    \alert{a register that does not read back is an unclocked peripheral}.}
\end{frame}

\begin{frame}{A Worked Investigation: The LED That Does Not Blink}{Six causes, ruled out in order}
    \begin{center}
    {\scriptsize
    \striped
    \begin{tabular}{@{}c l R{.58\textwidth}@{}}
        \toprule
        \textbf{Step} & \textbf{Question} & \textbf{The observation that settles it} \\
        \midrule
        1 & Is the clock enabled? & Write \texttt{0xFFFF} to \texttt{GPIOA->MODER} and read it back. If it reads zero, the \texttt{RCC->AHB1ENR} bit is missing. \\
        2 & Is the pin an output? & Read \texttt{MODER} bits 11:10. If not \texttt{01}, the driver is disconnected. \\
        3 & Is \texttt{ODR} changing? & Watch \texttt{ODR} bit 5 while stepping. If it never changes, the write is not running, or \texttt{volatile} is missing. \\
        4 & Does the pin move? & Meter or scope on PA5. \texttt{ODR} toggles but the pin does not: an alternate function or analog mode has claimed it. \\
        5 & Is it the pin you think? & Check the schematic. On some Nucleo variants the LED is not on PA5. \\
        6 & Is the circuit right? & Voltage but no light: a reversed LED, an open resistor, or the wrong sink/source direction. \\
        \bottomrule
    \end{tabular}}
    \end{center}

    \vspace{1mm}
    {\small The walk goes \alert{up the five levels}: register, configuration, C, pin voltage, circuit. Steps 1--4 use the debugger and a meter, not the source code. Rereading the program is almost always the slowest move.}
\end{frame}


%===============================================================================
\section{What You Have Learned}
%===============================================================================

\begin{frame}{What You Have Learned}
    \begin{columns}[T]
    \begin{column}{.48\textwidth}
        \textbf{The mechanism}
        {\small
        \begin{itemize}\itemsep1pt
            \item A CPU fetches, decodes, executes --- forever
            \item Memory is numbered; some numbers are hardware
            \item C reaches them through \alert{pointers}
            \item Instructions are what your C becomes
        \end{itemize}}
    \end{column}
    \begin{column}{.48\textwidth}
        \textbf{The practice}
        {\small
        \begin{itemize}\itemsep1pt
            \item Enable the clock first
            \item \texttt{volatile} on everything hardware touches
            \item Read--modify--write preserves the neighbours
            \item Interrupts for deadlines, timers for time
            \item \alert{Predict, then measure}
        \end{itemize}}
    \end{column}
    \end{columns}

    \vspace{3mm}
    \begin{block}{The one sentence from Chapter 1}
        \alert{The CPU has no special instruction for ``turn on an LED''.} It has
        instructions that read and write addresses --- and some of those addresses
        are wired to hardware.

        \vspace{1mm}
        Everything in the thirteen chapters since has been an elaboration of that.
    \end{block}
\end{frame}


%===============================================================================
\section{The Capstone Project}
%===============================================================================

\begin{frame}{The Capstone Project}{A system you design before you build}
    \leavevmode
    {\small Six weeks, in a group, on hardware you keep. It replaces a final exam question with a final exam \emph{system}: one you can plug in, measure, and defend.}

    \vspace{2mm}
    \begin{columns}[T]
    \begin{column}{.46\textwidth}
        \begin{block}{A shared core, then a free choice}
            {\small \textbf{Weeks 1--4:} everyone builds the same core: GPIO with interrupt-driven input, and a UART debug link, under the non-blocking architecture.

            \vspace{1mm}
            \textbf{Weeks 5--6:} each group adds one enhancement peripheral, integrates it without breaking the core, and defends it.}
        \end{block}
    \end{column}
    \begin{column}{.5\textwidth}
        {\footnotesize
        \striped
        \begin{tabular}{@{}c R{.62\textwidth}@{}}
            \toprule
            \textbf{Week} & \textbf{Deliverable} \\
            \midrule
            1 & M1: GPIO output, verified with a meter or scope \\
            2 & M2: debounced EXTI input drives the output through a \texttt{volatile} flag \\
            3 & M3: UART debug link at a known baud rate \\
            4 & M4: integration in one non-blocking loop \\
            5 & Enhancement chosen, brought up, integrated \\
            6 & Demonstration and written report \\
            \bottomrule
        \end{tabular}}
    \end{column}
    \end{columns}
\end{frame}

\begin{frame}{The Required Core}{By the end of Week 4, whatever the enhancement}
    \begin{enumerate}\itemsep3pt
        \item At least one \alert{GPIO output}, changed only through a register write the group can point to and explain.
        \item At least one external input handled through \alert{EXTI}, not by polling. The handler clears the flag, records the event, returns.
        \item A \alert{UART link} at a stated baud rate, reporting at least one state change to a PC terminal.
        \item The whole system as a \alert{non-blocking cooperative loop}: interrupts set flags, the main loop does the work, no branch blocks another.
        \item A one-paragraph \alert{timing budget}, produced \emph{before} Milestone 4 is demonstrated.
    \end{enumerate}

    \vspace{3mm}
    {\small Milestones 1--3 may use the peripherals separately. Milestone 4 needs them combined: one program, one main loop, one timing budget.}
\end{frame}

\begin{frame}{Milestone Rubrics (1)}{Each milestone is 10 points, weighted equally}
    \begin{columns}[T]
    \begin{column}{.48\textwidth}
        \textbf{M1: GPIO output}
        {\scriptsize
        \striped
        \begin{tabular}{@{}c R{.78\linewidth}@{}}
            \toprule
            \textbf{Pts} & \textbf{Criterion} \\
            \midrule
            4 & Pin configured (clock, \texttt{MODER}) and verified with a meter or scope, not just ``the LED lit''. \\
            3 & The group states from memory which bits were written and why. \\
            3 & Named masks, no magic numbers; clock enable before any write. \\
            \bottomrule
        \end{tabular}}
    \end{column}
    \begin{column}{.48\textwidth}
        \textbf{M2: Debounced EXTI input}
        {\scriptsize
        \striped
        \begin{tabular}{@{}c R{.78\linewidth}@{}}
            \toprule
            \textbf{Pts} & \textbf{Criterion} \\
            \midrule
            3 & Interrupt configured (\texttt{SYSCFG}, edge, \texttt{IMR}, NVIC) and reliable. \\
            3 & Debounce by timestamp, not a blocking delay in the handler. \\
            2 & Handler clears the pending bit first, sets a \texttt{volatile} flag, returns. \\
            2 & Main loop drives the M1 output correctly over at least ten presses. \\
            \bottomrule
        \end{tabular}}
    \end{column}
    \end{columns}
\end{frame}

\begin{frame}{Milestone Rubrics (2)}
    \begin{columns}[T]
    \begin{column}{.48\textwidth}
        \textbf{M3: UART debug link}
        {\scriptsize
        \striped
        \begin{tabular}{@{}c R{.78\linewidth}@{}}
            \toprule
            \textbf{Pts} & \textbf{Criterion} \\
            \midrule
            4 & Correct baud register value for the stated clock, verified by clean text on a PC terminal. \\
            3 & \texttt{TXE} polled correctly: no missed or duplicated characters. \\
            3 & The group states the baud rate, its clock, and the arithmetic, from memory. \\
            \bottomrule
        \end{tabular}}
    \end{column}
    \begin{column}{.48\textwidth}
        \textbf{M4: Integration}
        {\scriptsize
        \striped
        \begin{tabular}{@{}c R{.78\linewidth}@{}}
            \toprule
            \textbf{Pts} & \textbf{Criterion} \\
            \midrule
            4 & All three peripherals in one non-blocking loop, no branch blocking another. \\
            2 & Every ISR-shared variable is \texttt{volatile}, and the group says why. \\
            2 & A timing budget written \emph{before} the demonstration. \\
            2 & Live fault injection: the group finds the fault with the bottom-up method, not by guessing. \\
            \bottomrule
        \end{tabular}}
    \end{column}
    \end{columns}
\end{frame}

\begin{frame}{The Enhancement}{Weeks 5--6: one more peripheral, integrated}
    \begin{columns}[T]
    \begin{column}{.5\textwidth}
        {\small \textbf{Option A: DHT11 sensor.} One open-drain data line with a timing-based protocol. It needs bit-banging with microsecond delays, and a reading takes about 20\,ms. The group must decide: one blocking call scheduled as a rare low-priority task, or interrupt-driven phases. Justify it against the timing budget.

        \vspace{2mm}
        \textbf{Option B: WS2812 LEDs.} Each bit is a precisely timed pulse (about 1.25\,\textmu s), 24 bits per LED, no clock line. Generate it by cycle-counted bit-banging, or from a timer's PWM with DMA.

        \vspace{2mm}
        \emph{A group may propose another peripheral, with approval, if it brings a comparable integration challenge.}}
    \end{column}
    \begin{column}{.46\textwidth}
        \textbf{Enhancement rubric (20 points)}
        {\scriptsize
        \striped
        \begin{tabular}{@{}c R{.78\linewidth}@{}}
            \toprule
            \textbf{Pts} & \textbf{Criterion} \\
            \midrule
            6 & Works alone, verified independently (scope, reference reading). \\
            6 & Integrated without breaking M4; the timing budget is \alert{re-measured}, not assumed. \\
            4 & Integration timing cost estimated beforehand, checked afterward. \\
            4 & The group explains the hardest design decision, unprompted. \\
            \bottomrule
        \end{tabular}}
    \end{column}
    \end{columns}
\end{frame}

\begin{frame}{Final Report and Demonstration}{Week 6}
    \begin{columns}[T]
    \begin{column}{.48\textwidth}
        \begin{block}{Live demonstration}
            {\small Fault injection, as in Milestone 4. The instructor may disconnect a wire, ask for a register value, or ask what happens if an input arrives during the enhancement's busiest moment.}
        \end{block}
    \end{column}
    \begin{column}{.48\textwidth}
        \begin{block}{Written report: at most four pages}
            {\small
            \begin{itemize}\itemsep1pt
                \item functional decomposition
                \item final pin assignment
                \item timing budget for the complete system
                \item verification plan: what was measured, with what instrument, and the pass/fail criterion
            \end{itemize}}
        \end{block}
    \end{column}
    \end{columns}

    \vspace{3mm}
    \alert{A demonstration with no timing budget, and a timing budget with no demonstration, both earn partial credit at best.} The two claims are not substitutes for one another.
\end{frame}


%===============================================================================
\section{Summary}
%===============================================================================

\begin{frame}{Summary}
	\begin{columns}[c]
		\begin{column}{.62\textwidth}
			\begin{center}
				\begin{tikzpicture}[x=10mm, y=7.8mm]
					\node[blk, minimum width=56mm, minimum height=5.0mm, inner sep=1.2pt, font=\scriptsize] (n0) at (0,0.00) {Written requirement};
					\node[blk, minimum width=56mm, minimum height=5.0mm, inner sep=1.2pt, font=\scriptsize] (n1) at (0,-1.05) {Functional decomposition};
					\node[blk, minimum width=56mm, minimum height=5.0mm, inner sep=1.2pt, font=\scriptsize] (n2) at (0,-2.10) {Peripheral and pin selection};
					\node[blk, minimum width=56mm, minimum height=5.0mm, inner sep=1.2pt, font=\scriptsize] (n3) at (0,-3.15) {Polling, timed polling or interrupts};
					\node[blk, minimum width=56mm, minimum height=5.0mm, inner sep=1.2pt, font=\scriptsize] (n4) at (0,-4.20) {Timing budget before any code};
					\node[blk, minimum width=56mm, minimum height=5.0mm, inner sep=1.2pt, font=\scriptsize] (n5) at (0,-5.25) {Cooperative loop; protect shared data};
					\node[blk, minimum width=56mm, minimum height=5.0mm, inner sep=1.2pt, font=\scriptsize] (n6) at (0,-6.30) {Measure and compare with the budget};
					\node[cpublk, minimum width=56mm, minimum height=5.0mm, inner sep=1.2pt, font=\scriptsize, align=center] (n7) at (0,-7.65) {Predict first,\\then measure};
					\foreach \a/\b in {0/1,1/2,2/3,3/4,4/5,5/6,6/7}{\draw[flow] (n\a) -- (n\b);}
				\end{tikzpicture}
			\end{center}
		\end{column}
		\begin{column}{.34\textwidth}
			{\scriptsize
				\textbf{Terminology covered}\\[1mm]
				\begin{itemize}
					\item Functional decomposition
					\item Timing budget
					\item Cooperative main loop
					\item Shared data protection
					\item Debugging checklist
				\end{itemize}
				\vspace{3mm}
				\textbf{Last chapter}\\[1mm]
				Use this method in every design you build.}
		\end{column}
	\end{columns}
\end{frame}

\end{document}

